\begin{proposition}[The genus bound and delta-one curve germs]\hypertarget{tag-01B4}{}\label{cdp:st:A.5.14}
\StatementRecord{cdp:st:A.5.14}{proved}{01B4}{\SeriesRef{cdp:eq:A.J.31}, \SeriesRef{cdp:st:A.5.5}}

Let $f:Y\to C_0$ be a proper connected-fibre map from a
smooth compact complex threefold with $H^2(Y,\Q)=0$ to
a smooth curve. Suppose its smooth general fibre has
trivial canonical bundle. Let $mS$ be a full fibre with
irreducible reduced nonnormal support, and let $D$ be
a reduced irreducible component of its conductor image.
Then
\begin{equation}
                         p_a(D)\le2.                \hypertarget{tag-005Z}{}\label{cdp:eq:A.J.30}
\end{equation}
If the normalization of $D$ has positive genus, then
$D$ is a local complete intersection in $Y$; every
singular germ of $D$ is an ordinary node or an ordinary
cusp, and the sum of its delta invariants is at most
one. In particular, when the generic transverse type
along such a positive-genus $D$ is $A_2$, $A_3$, or an
ordinary triple point, $D$ is smooth elliptic and the
transverse analytic type is constant along all of $D$.
\end{proposition}


\begin{proof}

\begin{enumerate}[label={$\langle1\rangle$\arabic*.},leftmargin=*]
\item Restrict $f$ to a disc about the image of $mS$ and trivialize the
canonical line of the disc. The sheaf $f_*K_Y$ is torsion-free of rank
one, hence locally free there. The evaluation of a generator is
nonzero on the smooth nearby fibres, where their canonical lines are
trivial. Its zero divisor is therefore $aS$ for an integer $a\ge0$.
Thus $K_Y\simeq\OO_Y(aS)$ on this neighbourhood. The full-fibre equation
trivializes $\OO_Y(mS)$, and adjunction gives
$\omega_S\simeq\OO_Y((a+1)S)|_S$, a torsion line.
The Cartier sequence and Riemann--Roch on $Y$ give $\chi(\OO_S)=0$,
since $c_1(\OO_Y(S))=0$ in rational cohomology. A nonzero section of a
torsion line on the compact irreducible surface has a nonzero constant
tensor power; it therefore trivializes the line. Consequently
$e=h^0(S,\omega_S)\in\{0,1\}$, and Serre duality gives
$h^1(\OO_S)=1+e$. For $e=1$, the line $\omega_S$ is trivial.

\item Write $\eta:T\to S$ for normalization,
$\mathcal Q=\eta_*\OO_T/\OO_S$, and $\Gamma$ for
the full conductor image scheme. Equality of the
constant global functions on $S$ and $T$ gives the
injection $H^0(\mathcal Q)\hookrightarrow
H^1(\OO_S)$. To identify the conductor-duality map, put
$\mathcal B=\eta_*\OO_T$ and $\mathfrak c=\operatorname{Ann}(\mathcal Q)$.
The normal surface $T$ is Cohen--Macaulay. Thus $\mathcal B$ is a
maximal Cohen--Macaulay module on the Gorenstein surface $S$.
Multiplication in the total quotient sheaf identifies
$\mathcal Hom(\mathcal B,\omega_S)$ with $\mathfrak c\omega_S$.
Applying $\mathcal Hom(-,\omega_S)$ to
$0\to\OO_S\to\mathcal B\to\mathcal Q\to0$ therefore gives
$\mathcal Ext^1(\mathcal Q,\omega_S)=\omega_S|_\Gamma$ and vanishing
of the other sheaf Ext groups. The local-to-global Ext sequence and
Serre duality identify
$H^0(S,\mathcal Q)^*$ with
$H^1(\Gamma,\omega_S|_\Gamma)$. Taking dimensions gives
\begin{equation}
 h^1(\Gamma,\omega_S|_\Gamma)
       =h^0(S,\mathcal Q)\le1+e.                  \hypertarget{tag-0060}{}\label{cdp:eq:A.J.31}
\end{equation}
The restriction quotient
$\omega_S|_\Gamma\twoheadrightarrow\omega_S|_D$
has a kernel supported in dimension at most one, so
its $H^1$ map is surjective. For $e=1$ it follows that
$p_a(D)=h^1(\OO_D)\le2$. For $e=0$, put
$L=\omega_S|_D$. This torsion line has degree zero,
and \eqref{cdp:eq:A.J.31} gives $h^1(D,L)\le1$. Curve
Riemann--Roch gives
\[
 p_a(D)=1+h^1(D,L)-h^0(D,L)\le2,
\]
proving \eqref{cdp:eq:A.J.30} in both cases.

\item Let $\nu:\widehat D\to D$ be normalization. The
normalization sequence gives the identity
\begin{equation}
 p_a(D)=g(\widehat D)+
        \sum_{x\in D}\dim_{\C}
          (\nu_*\OO_{\widehat D}/\OO_D)_x.       \hypertarget{tag-0061}{}\label{cdp:eq:A.J.32}
\end{equation}
If $g(\widehat D)>0$, every local delta is consequently
at most one and their sum is at most one. We give the
analytic local rings at a point with delta one.
Let $A=\OO_{D,x}$ and let
$B=\bigoplus_{i=1}^r\C\{t_i\}$ be its finite
normalization. Here $\dim_{\C}(B/A)=1$. The residue map
of the local ring $A$ has the same value on every
normalization branch. The resulting quotient of
$\C^r$ by its diagonal has dimension $r-1$, so
$r\le2$.

\item For $r=2$, the subring
\[
 A_0=\{(b_1,b_2)\in\C\{t_1\}\oplus\C\{t_2\}:
                                  b_1(0)=b_2(0)\}
\]
contains $A$ and has colength one in $B$. Hence
$A=A_0$. The explicit homomorphism
$\C\{u,v\}/(uv)\to A_0$ sends
$u$ to $(t_1,0)$, $v$ to $(0,t_2)$, and a constant
to the corresponding equal pair. Expanding a
convergent series shows that it is an isomorphism.
This is the ordinary node germ.

\item For $r=1$, an element of $A$ with valuation one would
be a local coordinate on the normalization disc by
the holomorphic inverse function theorem. Convergent
composition in the analytic algebra $A$ would then
give $\C\{t\}\subset A$, contrary to its colength
one. Therefore every nonconstant element of $A$ has
valuation at least two. We have
$A\subset\C+t^2\C\{t\}$; both rings have
colength one in $\C\{t\}$, so equality holds.
The map
\[
 \C\{u,v\}/(v^2-u^3)\longrightarrow A,
                 \qquad (u,v)\longmapsto(t^2,t^3)
\]
is an isomorphism: every series with no linear term
is uniquely the sum of an even series in $t^2$ and
$t^3$ times a series in $t^2$. This is the ordinary
cusp germ. A delta-zero germ is its smooth
normalization disc.

\item We finally check the embedding into the threefold $Y$. These analytic rings have embedding
dimension at most two. In local coordinates of the
smooth threefold, the tangent space of $D$ therefore
has codimension at least one. There is an element
$z$ of its ideal with $dz(x)\ne0$. The inverse
function theorem makes $z$ an ambient coordinate,
and the ideal of $D$ is the inverse image of a
height-one ideal in the regular analytic ring
$\C\{u,v\}=\OO_{Y,x}/(z)$. This ideal is principal:
a pure reduced curve in a smooth surface is a Cartier
divisor, or equivalently its reduced height-one ideal
is a product of the prime factors in this factorial
ring. Thus $I_{D,x}=(z,h)$, a regular sequence.
This proves the asserted lci embedding, including
at the node and cusp. Proposition~\ref{cdp:st:A.5.5} applies
to the three specified generic transverse types
and proves the final smooth elliptic and constant-type
conclusion.

\end{enumerate}
\end{proof}